\scanpage{73}

ayant donc la base

\[1,x,y,x^2,xy,y^2\]

les valeurs propres de $u$ correspondantes sont resp\textsuperscript{t}.

\[1,\zeta,\zeta^{-1},\unsure{\zeta|\zeta^2},1,\zeta^{-2}\]

donc les valeurs propres qui y figurent sont \deleted{\unsure{respectivement}} parmi les \inserted{$5$} valeurs

\[\zeta^{-2},\zeta^{-1},\zeta^0=1,\zeta^1=\zeta,\zeta^2.\]

Celles-ci sont \deleted{\unsure{ttes|toutes}} distinctes \unsure{ssi|sss|si} $n\geq5$, dans ce cas les \deleted{fonctions}\inserted{poly} quasi-invariantes de degré $\leq2$ \inserted{sous $G^+$} sont exact\textsuperscript{t} les fonctions

\[\begin{aligned}
a+bxy&&\text{valeur propre }&\zeta^0=1\\
ax\quad(\text{resp. }ay)&&\text{valeur propre }&\zeta\quad(\text{resp }\zeta^{-1})\\
ax^2\quad(\text{resp }ay^2)&&\text{valeur propre }&\zeta^2\quad(\text{resp }\zeta^{\unsure{-1|-2}})
\end{aligned}\]

Si on veut de plus la quasi-invariance sous $G\simeq D_n$ tout entier, \deleted{sachant que $\sigma$}\inserted{i.e. sous $\sigma$} notant que $\sigma(ax)=ay$, $\sigma(ay)=ax$, $\sigma(ax^2)=ay^{2}$, $\sigma(ay^2)=ax^2$, on trouve que les seules fonctions \inserted{poly. de degré $\leq2$} quasi-invariantes sous $G$ sont \deleted{\unsure{les}} $a+bxy$ (qui sont \unsure{$=$|en|{}} invariantes)

Donc les seules côniques invariantes sous $G$ sont les coniques d'équation $a+bxy$ ; donc si, on veut les coniques circonscrites à $\Pi_\zeta$ il faut $a+b=0$ i.e. $b=-a$, on trouve la \deleted{\unsure{seule}} cônique \uline{invariante} circonscrite, d'équation $xy-1=0$ ....

% archive sheet number: 72
